\documentclass[a4paper,12pt]{article}

\usepackage[T2A]{fontenc}
\usepackage[utf8]{inputenc}
\usepackage[russian]{babel}
\usepackage{amsmath,amssymb}
\usepackage{geometry}

\geometry{left=2cm,right=2cm,top=2cm,bottom=2cm}
\setlength{\parindent}{0pt}
\setlength{\parskip}{0.65em}

\begin{document}

\begin{titlepage}
    \centering
    \vspace*{0.25\textheight}
    {\LARGE Расчетное задание 2.\par}
    \vspace{1.5cm}
    {\Large Шубин Егор Вячеславович\par}
    {\Large P3209\par}
    \vfill
    {\large 2026\par}
\end{titlepage}

\section*{Задача 1}

Для дуги окружности радиуса \(r\), видимой из точки \(O\) под углом
\(\theta\), поле в центре равно
\[
    B_{\text{дуги}}=\frac{\mu_0 I}{4\pi r}\theta .
\]

а) Для данного рисунка дуга окружности радиуса \(a\) видна под углом \(2\pi-\varphi\), а дуга радиуса \(b\) видна под углом \(\varphi\). Поэтому
\[
    B=\frac{\mu_0 I}{4\pi}
    \left(\frac{2\pi-\varphi}{a}+\frac{\varphi}{b}\right).
\]

б) Круговая часть контура имеет угол \(3\pi/2\), а прямая часть контура видна под углом \(45^\circ\). Поэтому
\[
    B_{\text{окр}}=\frac{\mu_0 I}{4\pi}\frac{3\pi}{2a},
    \qquad
    B_{\text{пр}}=
    2\cdot\frac{\mu_0 I}{4\pi b}\sin45^\circ
    =\frac{\mu_0 I}{4\pi}\frac{\sqrt2}{b}.
\]
Следовательно,
\[
    B=\frac{\mu_0 I}{4\pi}
    \left(\frac{3\pi}{2a}+\frac{\sqrt2}{b}\right)
\]
Поле направлено внутрь листа.

\section*{Задача 2}

Выделим на сфере кольцо между углами \(\theta\) и \(\theta+d\theta\):
\[
    dq=\sigma\cdot 2\pi R^2\sin\theta\,d\theta,
    \qquad
    dI=\frac{\omega\,dq}{2\pi}
    =\sigma\omega R^2\sin\theta\,d\theta .
\]
Поле этого кольца в центре сферы:
\[
    dB=\frac{\mu_0 dI(R\sin\theta)^2}{2R^3}
    =\frac{\mu_0\sigma\omega R}{2}\sin^3\theta\,d\theta .
\]
Тогда
\[
    B=\frac{\mu_0\sigma\omega R}{2}\int_0^\pi\sin^3\theta\,d\theta
    =\frac{2}{3}\mu_0\sigma\omega R.
\]
\[
    B=\frac{2}{3}\cdot 4\pi\cdot10^{-7}\cdot10^{-5}\cdot70\cdot0.05
    \approx 2.93\cdot10^{-11}\ \text{Тл}.
\]
\[
  B\approx 2.93\cdot10^{-11}\ \text{Тл}
\]

\section*{Задача 3}

Для одного витка сила вдоль оси равна \(dF=I\,d\Phi\). На участке длины
\(dz\) находится \(n\,dz\) витков, поэтому
\[
    F_z=nI\int d\Phi=nI(\Phi_2-\Phi_1).
\]
\[
    F=nI|\Phi_2-\Phi_1|
\]

\section*{Задача 4}

Поверхностная плотность тока:
\[
    K=\frac{I}{2\pi R}.
\]
Внутри цилиндра \(B_{\text{вн}}=0\), снаружи у поверхности
\[
    B_{\text{нар}}=\frac{\mu_0 I}{2\pi R}=\mu_0 K.
\]
\[
    p=\frac{B_{\text{нар}}^2-B_{\text{вн}}^2}{2\mu_0}
    =\frac{\mu_0 I^2}{8\pi^2R^2}.
\]
\[
    p=\frac{4\pi\cdot10^{-7}\cdot 50^2}
    {8\pi^2\cdot(0.05)^2}
    \approx 1.59\cdot10^{-2}\ \text{Па}.
\]


\section*{Задача 5}

\[
    B(t)=\dot B\,t,\qquad
    x(t)=\frac{at^2}{2},\qquad
    S(t)=lx(t)=\frac{lat^2}{2}.
\]
Поток:
\[
    \Phi(t)=B(t)S(t)=\frac{\dot B\,l a t^3}{2}.
\]
ЭДС индукции:
\[
    \mathcal{E}=\left|\frac{d\Phi}{dt}\right|
    =\frac{3}{2}\dot B\,l a t^2.
\]
\[
    \mathcal{E}=
    \frac{3}{2}\cdot0.1\cdot0.20\cdot0.10\cdot2^2
    =1.2\cdot10^{-2}\ \text{В}.
\]

\section*{Задача 6}

Потокосцепление тора:
\[
    \Psi=N\Phi.
\]
Индуктивность:
\[
    L=\frac{\Psi}{I}=\frac{N\Phi}{I}.
\]
Энергия магнитного поля:
\[
    W=\frac{LI^2}{2}=\frac{N\Phi I}{2}.
\]
\[
    W=\frac{500\cdot10^{-3}\cdot2}{2}=0.50\ \text{Дж}.
\]

\section*{Задача 7}

Скорость протона после ускорения:
\[
    eU=\frac{m_pv^2}{2}.
\]
Радиус траектории в магнитном поле:
\[
    r=\frac{m_pv}{eB}
    =\frac{1}{B}\sqrt{\frac{2m_pU}{e}}.
\]
\[
    r=\frac{1}{0.51}\sqrt{\frac{2\cdot1.67\cdot10^{-27}\cdot5\cdot10^5}
    {1.60\cdot10^{-19}}}
    \approx 0.2\ \text{м}.
\]
Так как толщина области поля \(d\) отсчитывается вдоль первоначального
направления движения,
\[
    \sin\alpha=\frac{d}{r}=\frac{0.1}{0.2}\approx 0.5.
\]
Следовательно,
\[
    \alpha= 30.0^\circ
\]

\section*{Задача 8}

Амплитуда тока в контуре убывает по закону
\[
    I_m(t)=I_{m0}e^{-\beta t},
    \qquad
    \beta=\frac{R}{2L}.
\]
\[
    t_e=\frac{1}{\beta}=\frac{2L}{R}.
\]
\[
    T=\frac{2\pi}{\omega},
    \qquad
    \omega=\sqrt{\frac{1}{LC}-\frac{R^2}{4L^2}}.
\]
\[
    N=\frac{t_e}{T}=\frac{L\omega}{\pi R}.
\]
При \(C=10\,\text{мкФ}\), \(L=25\,\text{мГн}\), \(R=1\,\text{Ом}\):
\[
    \omega=\sqrt{\frac{1}{25\cdot10^{-3}\cdot10\cdot10^{-6}}
    -\frac{1}{4(25\cdot10^{-3})^2}}
    \approx 2.00\cdot10^3\ \text{с}^{-1},
\]
\[
    N=\frac{25\cdot10^{-3}\cdot2.00\cdot10^3}{\pi}
    \approx 15.9\ \text{колебаний}.
\]

\end{document}
